% ╔══════════════════════════════════════════════════════════════════════════════╗
% ║  PHYS4071 · Astrophysics & Cosmology — Course Notes                        ║
% ║  Compile with: pdflatex → bibtex → pdflatex → pdflatex                     ║
% ╚══════════════════════════════════════════════════════════════════════════════╝

\documentclass[12pt, a4paper]{article}

% ─── Encoding & Fonts ──────────────────────────────────────────────────────────
\usepackage[T1]{fontenc}
\usepackage[utf8]{inputenc}
\usepackage{mathpazo}                  % Palatino for text + math
\usepackage{microtype}                 % Microtypography (better justification)

% ─── Page Layout ───────────────────────────────────────────────────────────────
\usepackage[
  a4paper,
  top=2.8cm, bottom=2.8cm,
  left=2.8cm, right=2.8cm,
  headheight=15pt
]{geometry}

% ─── Mathematics ───────────────────────────────────────────────────────────────
\usepackage{amsmath, amssymb, amsfonts}
\usepackage{bm}                        % Bold math symbols

% ─── Color Palette  ────────────────────────────────────────────────────────────
%     "Midnight Observatory" — navy anchors, gold & teal accents, clean white page
\usepackage{xcolor}
\definecolor{DeepSpace}{HTML}{0D1B2A}       % very dark navy   – headings
\definecolor{StarGold}{HTML}{C9A84C}        % warm gold        – intuition
\definecolor{NebulaTeal}{HTML}{2E86AB}      % teal             – theory
\definecolor{ForestGreen}{HTML}{2D6A4F}     % deep green       – example
\definecolor{SlateGray}{HTML}{617A8A}       % cool slate       – extra
\definecolor{GoldBG}{HTML}{FDF8EC}          % pale gold bg
\definecolor{TealBG}{HTML}{EBF5FA}          % pale teal bg
\definecolor{GreenBG}{HTML}{EBF5EF}         % pale green bg
\definecolor{GrayBG}{HTML}{F3F4F6}          % pale gray bg
\definecolor{TitleBG}{HTML}{090F1A}         % near-black title page

% ─── Graphics & TikZ ───────────────────────────────────────────────────────────
\usepackage{tikz}
\usetikzlibrary{calc, backgrounds}
\usepackage{graphicx}

% ─── Boxes (tcolorbox) ─────────────────────────────────────────────────────────
\usepackage[most, breakable]{tcolorbox}

%   ── Physics Intuition ──────────────────────────────────────────────────────
\newtcolorbox{intuition}[1][]{
  breakable, enhanced,
  colback   = GoldBG,
  colframe  = StarGold,
  leftrule  = 4pt, rightrule = 0pt, toprule = 0pt, bottomrule = 0pt,
  arc = 0pt, outer arc = 0pt,
  before upper = {%
    \textcolor{StarGold}{\textbf{\small\scshape Physics Intuition}}%
    \par\smallskip},
  top = 2pt, bottom = 6pt, left = 8pt, right = 8pt,
  #1
}

%   ── Theory ─────────────────────────────────────────────────────────────────
\newtcolorbox{theory}[1][]{
  breakable, enhanced,
  colback   = TealBG,
  colframe  = NebulaTeal,
  leftrule  = 4pt, rightrule = 0pt, toprule = 0pt, bottomrule = 0pt,
  arc = 0pt, outer arc = 0pt,
  before upper = {%
    \textcolor{NebulaTeal}{\textbf{\small\scshape Theory}}%
    \par\smallskip},
  top = 2pt, bottom = 6pt, left = 8pt, right = 8pt,
  #1
}

%   ── Question Example (auto-numbered per section) ───────────────────────────
\newtcbtheorem[number within=section]{example}{Question Example}{
  breakable, enhanced,
  colback    = GreenBG,
  colframe   = ForestGreen,
  arc = 3pt, outer arc = 3pt,
  leftrule   = 1.2pt, rightrule = 1.2pt,
  toprule    = 1.2pt, bottomrule = 1.2pt,
  fonttitle  = \bfseries\small\color{ForestGreen},
  coltitle   = ForestGreen,
  attach boxed title to top left = {yshift=-2mm, xshift=8pt},
  boxed title style = {
    colback = white, colframe = ForestGreen,
    arc = 2pt, outer arc = 2pt,
    top = 1pt, bottom = 1pt, left = 4pt, right = 4pt
  },
  top = 6pt, bottom = 6pt, left = 8pt, right = 8pt,
}{ex}

%   ── Extra ──────────────────────────────────────────────────────────────────
\newtcolorbox{extra}[1][]{
  breakable, enhanced,
  colback   = GrayBG,
  colframe  = SlateGray,
  leftrule  = 2pt, rightrule = 2pt, toprule = 0.6pt, bottomrule = 0.6pt,
  arc = 3pt, outer arc = 3pt,
  before upper = {%
    \textcolor{SlateGray}{\textbf{\small\scshape Extra}}%
    \par\smallskip},
  top = 2pt, bottom = 6pt, left = 8pt, right = 8pt,
  #1
}

% ─── Section Title Style ───────────────────────────────────────────────────────
\usepackage{titlesec}
\titleformat{\section}[block]
  {\large\bfseries\color{DeepSpace}}
  {\thesection.}{0.5em}{}
  [\vspace{2pt}\color{StarGold}\rule{\textwidth}{0.7pt}]
\titlespacing{\section}{0pt}{1.6em}{0.9em}

\titleformat{\subsection}[block]
  {\normalsize\bfseries\color{DeepSpace}}
  {\thesubsection}{0.6em}{}
  [\vspace{2pt}\color{NebulaTeal!60}\rule{\textwidth}{0.4pt}]
\titlespacing{\subsection}{0pt}{1.2em}{0.5em}

% ─── Header / Footer ───────────────────────────────────────────────────────────
\usepackage{fancyhdr}
\pagestyle{fancy}
\fancyhf{}
\renewcommand{\headrulewidth}{0pt}
\fancyhead[L]{\small\color{SlateGray}\itshape PHYS4071 \,·\, Astrophysics \& Cosmology}
\fancyhead[R]{\small\color{SlateGray}\thepage}
\fancyfoot[C]{\color{SlateGray!40}\rule{0.35\textwidth}{0.4pt}}

% ─── Hyperlinks ────────────────────────────────────────────────────────────────
\usepackage[hidelinks]{hyperref}
\hypersetup{
  colorlinks = true,
  linkcolor  = NebulaTeal,
  citecolor  = NebulaTeal,
  urlcolor   = NebulaTeal,
  pdftitle   = {PHYS4071 Cosmology Notes},
}

% ─── Miscellaneous ─────────────────────────────────────────────────────────────
\usepackage{enumitem}
\setlist[itemize]{leftmargin=*, label={\small\textcolor{StarGold}{$\bullet$}}}
\setlist[enumerate]{leftmargin=*}
\usepackage{booktabs}
\usepackage{lipsum}     % placeholder — remove once real content is added

% ─── Convenience macros ────────────────────────────────────────────────────────
\newcommand{\Hubble}{H_0}
\newcommand{\Mpc}{\,\mathrm{Mpc}}
\newcommand{\Gyr}{\,\mathrm{Gyr}}
\newcommand{\kms}{\,\mathrm{km\,s^{-1}}}

% ══════════════════════════════════════════════════════════════════════════════
\begin{document}
\pagestyle{empty}

% ─── Title Page ────────────────────────────────────────────────────────────────
\begin{tikzpicture}[remember picture, overlay]
  % Solid dark background
  \fill[TitleBG] (current page.south west) rectangle (current page.north east);

  % Starfield — random white dots of varied size and opacity
  \foreach \i in {1,...,220}{%
    \pgfmathsetseed{\i * 71 + 13}%
    \pgfmathsetmacro{\sx}{rand * 10.5 + 10.5}%
    \pgfmathsetmacro{\sy}{rand * 14.85 + 14.85}%
    \pgfmathsetmacro{\ro}{0.15 + 0.55 * rnd}%
    \pgfmathsetmacro{\rr}{0.4 + 1.1 * rnd}%
    \fill[white, opacity=\ro] (\sx cm, \sy cm) circle (\rr pt);%
  }%

  % Constellation sketch (Orion belt region, schematic)
  \newcommand{\starpt}[2]{\fill[StarGold, opacity=0.75] (#1 cm, #2 cm) circle (2.2pt);}%
  \starpt{4.0}{6.0}   \starpt{5.5}{7.8}   \starpt{6.8}{7.0}
  \starpt{8.2}{8.5}   \starpt{9.0}{7.2}   \starpt{6.8}{5.2}
  \starpt{5.8}{4.2}
  \draw[StarGold, opacity=0.25, line width=0.6pt]
    (4.0cm,6.0cm) -- (5.5cm,7.8cm) -- (6.8cm,7.0cm) --
    (8.2cm,8.5cm)  (6.8cm,7.0cm) -- (9.0cm,7.2cm)
    (6.8cm,7.0cm) -- (6.8cm,5.2cm) -- (5.8cm,4.2cm);

  % Horizontal gold accent rule
  \draw[StarGold, line width=1.4pt]
    ([shift={(2.0cm,-10.2cm)}]current page.north west) --
    ([shift={(-2.0cm,-10.2cm)}]current page.north east);
  \draw[StarGold!35, line width=0.5pt]
    ([shift={(2.0cm,-10.42cm)}]current page.north west) --
    ([shift={(-2.0cm,-10.42cm)}]current page.north east);
\end{tikzpicture}

\null\vspace{5.8cm}

\begin{center}
  {\small\color{StarGold}\textsc{\large P\kern1pt H\kern1pt Y\kern1pt S\kern2pt 4\kern1pt 0\kern1pt 7\kern1pt 1}}\\[0.7em]
  {\fontsize{32}{38}\selectfont\bfseries\color{white}Astrophysics \&}\\[0.15em]
  {\fontsize{32}{38}\selectfont\bfseries\color{white}Cosmology}\\[1.0em]
  {\color{StarGold!55}\rule{5.5cm}{0.5pt}}\\[0.7em]
  {\color{white!65}\normalsize Course Notes}\\[0.3em]
  {\color{white!40}\small Spring 2026}
\end{center}

\clearpage
\pagestyle{fancy}

% ─── Table of Contents ─────────────────────────────────────────────────────────
{%
  \hypersetup{linkcolor=DeepSpace}%
  \tableofcontents
}
\clearpage


% ══════════════════════════════════════════════════════════════════════════════
%  §1  THE EXPANDING UNIVERSE
% ══════════════════════════════════════════════════════════════════════════════
\section{The Expanding Universe}

% ── Physics Intuition ──────────────────────────────────────────────────────────
\begin{intuition}
\lipsum[1][1-4]   %TODO: replace — motivation / historical background / big-picture intuition
\end{intuition}

\bigskip

% ── Theory ─────────────────────────────────────────────────────────────────────
\begin{theory}
\lipsum[2][1-3]   %TODO: replace — derivation / assumptions / connection to intuition

The \textbf{Hubble–Lemaître Law} relates the recession velocity of a galaxy to its
proper distance~$d$:
\begin{equation}
  v \;=\; H_0 \, d,
  \label{eq:hubble}
\end{equation}
where $H_0 \approx 67.4 \kms\Mpc^{-1}$ is the Hubble constant (Planck~2018).
The cosmological redshift is defined via the scale factor $a(t)$:
\begin{equation}
  1 + z \;=\; \frac{a(t_0)}{a(t_e)} \;=\; \frac{1}{a(t_e)},
  \label{eq:redshift}
\end{equation}
using the convention $a(t_0)=1$ today.

\lipsum[2][4-5]   %TODO: replace
\end{theory}

\bigskip

% ── Question Example ───────────────────────────────────────────────────────────
\begin{example}{Recession Velocity from Redshift}{}
\lipsum[3][1-2]   %TODO: replace — question statement

\medskip\textbf{Solution.}\quad
\lipsum[3][3-4]   %TODO: replace — solution steps and answer
\end{example}

\bigskip

% ── Extra ──────────────────────────────────────────────────────────────────────
\begin{extra}
\lipsum[4][1-2]   %TODO: replace — aside, fun fact, or deeper reference
\end{extra}

\clearpage


% ══════════════════════════════════════════════════════════════════════════════
%  §2  THE FRIEDMANN EQUATIONS
% ══════════════════════════════════════════════════════════════════════════════
\section{The Friedmann Equations}

% ── Physics Intuition ──────────────────────────────────────────────────────────
\begin{intuition}
\lipsum[5][1-4]   %TODO: replace
\end{intuition}

\bigskip

% ── Theory ─────────────────────────────────────────────────────────────────────
\begin{theory}
The \textbf{FLRW metric} describes a homogeneous, isotropic universe:
\begin{equation}
  ds^2 = -c^2\,dt^2 + a(t)^2
    \!\left[\frac{dr^2}{1-kr^2} + r^2\,d\theta^2 + r^2\!\sin^2\!\theta\,d\phi^2\right],
\end{equation}
with curvature parameter $k\in\{-1,0,+1\}$.
Inserting into the Einstein field equations gives the \textbf{Friedmann equations}:
\begin{align}
  \left(\frac{\dot{a}}{a}\right)^{\!2}
    &= \frac{8\pi G}{3}\rho - \frac{kc^2}{a^2} + \frac{\Lambda c^2}{3},
    \label{eq:friedmann1}\\[6pt]
  \frac{\ddot{a}}{a}
    &= -\frac{4\pi G}{3}\!\left(\rho + \frac{3p}{c^2}\right) + \frac{\Lambda c^2}{3}.
    \label{eq:friedmann2}
\end{align}

\lipsum[6][1-4]   %TODO: replace — explain assumptions, equation of state $w=p/\rho c^2$
\end{theory}

\bigskip

% ── Question Example ───────────────────────────────────────────────────────────
\begin{example}{Scale Factor of a Flat Matter-Dominated Universe}{}
\lipsum[7][1-2]   %TODO: replace — question

\medskip\textbf{Solution.}\quad
For $k=0$, $\Lambda=0$ and pressureless dust ($p=0$), Eq.~\eqref{eq:friedmann1}
reduces to $\dot{a}^2 \propto a^{-1}$, giving:
\begin{equation}
  a(t) \;\propto\; t^{2/3}.
\end{equation}
\lipsum[7][3-4]   %TODO: replace — interpretation
\end{example}

\bigskip

% ── Extra ──────────────────────────────────────────────────────────────────────
\begin{extra}
\lipsum[8][1-2]   %TODO: replace — e.g.\ historical note on Friedmann (1922), Lemaître (1927)
\end{extra}

\clearpage


% ══════════════════════════════════════════════════════════════════════════════
%  §3  ENERGY BUDGET: THE ΛCDM MODEL
% ══════════════════════════════════════════════════════════════════════════════
\section[Energy Budget: LCDM Model]{Energy Budget: The $\Lambda$CDM Model}

% ── Physics Intuition ──────────────────────────────────────────────────────────
\begin{intuition}
\lipsum[9][1-4]   %TODO: replace — why dark matter + dark energy? what do observations say?
\end{intuition}

\bigskip

% ── Theory ─────────────────────────────────────────────────────────────────────
\begin{theory}
Define the \textbf{critical density} and dimensionless density parameters:
\begin{equation}
  \rho_\mathrm{crit}(z) = \frac{3H(z)^2}{8\pi G}, \qquad
  \Omega_i(z) = \frac{\rho_i(z)}{\rho_{\mathrm{crit}}(z)}.
\end{equation}
Each component scales with redshift according to its equation of state $w_i$:
\begin{equation}
  \rho_i(z) = \rho_{i,0}(1+z)^{3(1+w_i)}.
\end{equation}

In the \textbf{flat $\Lambda$CDM} model ($\Omega_k=0$,
$\Omega_\mathrm{m,0}+\Omega_{\Lambda}+\Omega_\mathrm{r,0}=1$):
\begin{equation}
  \frac{H(z)^2}{H_0^2}
    = \Omega_{\mathrm{m},0}(1+z)^3
    + \Omega_{\mathrm{r},0}(1+z)^4
    + \Omega_\Lambda.
  \label{eq:lcdm_hubble}
\end{equation}
Current Planck~2018 best-fit values:
$\Omega_\mathrm{m,0}\approx0.315$,
$\Omega_\Lambda\approx0.685$,
$\Omega_\mathrm{r,0}\approx9.2\times10^{-5}$.

\lipsum[10][1-3]   %TODO: replace
\end{theory}

\bigskip

% ── Question Example ───────────────────────────────────────────────────────────
\begin{example}{Age of the Universe in Flat $\Lambda$CDM}{}
\lipsum[11][1-2]   %TODO: replace

\medskip\textbf{Solution.}\quad
The lookback-time integral is:
\begin{equation}
  t_0 = \int_0^\infty \frac{dz}{(1+z)\,H(z)},
\end{equation}
with $H(z)$ from Eq.~\eqref{eq:lcdm_hubble}.
\lipsum[11][3-4]   %TODO: replace — numerical result and physical interpretation
\end{example}

\bigskip

% ── Extra ──────────────────────────────────────────────────────────────────────
\begin{extra}
\lipsum[12][1-2]   %TODO: replace — e.g.\ tension between $H_0$ measurements (CMB vs.\ local)
\end{extra}

% ══════════════════════════════════════════════════════════════════════════════
\end{document}
